% 天文学（综述）
% license CCBYSA3
% type Wiki

本文根据 CC-BY-SA 协议转载翻译自维基百科\href{https://en.wikipedia.org/wiki/Astronomy?utm_source=chatgpt.com}{相关文章}。

\begin{figure}[ht]
\centering
\includegraphics[width=6cm]{./figures/ddcd5c7b339e766a.png}
\caption{} \label{fig_Astron_1}
\end{figure}
天文学是一门研究天体及宇宙中发生的各种现象的自然科学。它运用数学、物理学和化学来解释这些天体的起源及其整体演化。研究对象包括行星、卫星、恒星、星云、星系、流星体、小行星和彗星等；相关的现象则包括超新星爆发、伽马射线暴、类星体、耀变体、脉冲星以及宇宙微波背景辐射等。更广义地说，天文学研究的一切事物都起源于地球大气层之外。而宇宙学（Cosmology）则是天文学的一个分支，专门研究宇宙整体的起源、结构与演化。

天文学是最古老的自然科学之一。在有文字记载的早期文明中，人类就已开始系统地观测夜空。其中包括古埃及人、巴比伦人、希腊人、印度人、中国人、玛雅人，以及美洲的许多原住民。在古代，天文学的范畴十分广泛，涵盖了天体测量学、天体导航、观测天文学以及历法制定等多个领域。

现代专业天文学大体上分为观测天文学与理论天文学两个分支。  
观测天文学的核心任务是通过望远镜与其他仪器对天体进行观测，从而获取数据；然后利用物理学的基本原理对这些数据进行分析。  
理论天文学则侧重于建立计算机模型或解析模型，用以描述天体及其物理现象。  
这两个分支相辅相成：理论天文学旨在解释观测结果，而观测天文学则通过实验数据来验证理论预测。

天文学是少数几个业余爱好者仍然能够发挥积极作用的科学领域之一。这一点在瞬变天体事件（如新星爆发、彗星出现等）的发现与观测中尤为明显。许多重要的天文发现——包括新的彗星——都是由业余天文学家首先发现的。


